\par\noindent\textbf{(1)} 连续性给 ∂\_tδρ+ρ₀∂\_xu=0，纵向动量为 ρ₀∂\_tu=−∂\_xδp+ζ\_eff∂\_x²u，局部绝热状态方程 δp=c²δρ。联立消去u得到 δρ\_tt−c²δρ\_xx−(ζ\_eff/ρ₀)δρ\_txx=0，耗散项必须有速度时间导数。
\par\noindent\textbf{(2)} 代入平面波得 ω²+i(ζ\_eff/ρ₀)k²ω−c²k²=0，令ν\_l=ζ\_eff/ρ₀，则 ω\_±=−iν\_lk²/2±√(c²k²−ν\_l²k⁴/4)。时间因子e\textasciicircum{}\{−iωt\}要求衰减对应Imω<0，所选方程满足这一物理符号。
\par\noindent\textbf{(3)} 弱黏性ν\_l k≪2c时，Reω≈±ck，Imω=−ν\_lk²/2，即声幅按exp(−ν\_lk²t/2)衰减。声压与速度同相到零阶，声波机械能随振幅平方而衰减更快；高阶相速度修正是O(ν\_l²k²/c²)。
\par\noindent\textbf{(4)} 吸收率随k²增长，说明短波声学对黏性更敏感；当ν\_lk>2c根号纯虚，本二阶连续介质模型形式上不再有振荡分支。极短波长时分子平均自由程和热扩散不再可忽略，不能把此过阻尼公式无限外推。
